\chapter*{\RL{ملخص}}
\addcontentsline{toc}{chapter}{Arabic Abstract}

\begin{RLtext}

\noindent يقدّم هذا التقرير مشروع نهاية الدراسات المنجز بشركة إس تي مايكروإلكترونيكس بتونس بهدف مقارنة منظومة \LR{STM32Cube} بمنظومات متحكمات دقيقة منافسة من خلال سيناريوهات برمجية مضبوطة. اعتمدنا وظائف متكافئة وشروط بناء موثقة، واستخرجنا البصمة الذاكرية من ملفات الربط، كما راجعنا التكافؤ الدلالي وعدالة المقارنة.

طوّرنا أداتين مساعدتين: \LR{MCU Workflow Studio} لتنظيم تحليل ملفات الربط بين المصنّعين، و\LR{MCU Benchmark Copilot Agent} لتنظيم إنشاء المشاريع ونقلها وبنائها وتتبع مصادرها والتحقق منها. تساعد الأداتان على توثيق سير العمل، لكن صحة كل نتيجة تبقى مرتبطة بأدلة البناء والقياس والعتاد.

غطّت مقارنة \LR{ST--GD32} سيناريوهات \LR{USB HID} و\LR{CDC} وثلاثة سيناريوهات لـ\LR{FreeRTOS}. استهلك متحكم إس تي ذاكرة فلاش أكثر بنسبة 35--88\%، واستخدم ذاكرة \LR{RAM} أقل بنسبة 57--65\% في سيناريوهات \LR{USB}. أظهرت القياسات المسجلة فرقاً قدره 61,7\% في زمن تهيئة \LR{USB} لصالح إس تي، غير أن غياب العينات المتكررة وحدود القياس المفصلة يمنع تعميم تفسير سببي.

جمعت مقارنة \LR{ST--Texas Instruments} بين تحليل ثابت للمشغلات واثني عشر مشروعاً متكافئاً. استخدم متحكم إس تي فلاشاً إجمالياً أكثر بنسبة 2,8\% وذاكرة \LR{RAM} أقل بنسبة 5,0\%. لا تسمح الأدلة الحالية باستنتاجات حول زمن التنفيذ أو الطاقة قبل توحيد الساعات وإعدادات الوحدات وحدود القياس وبروتوكول أخذ العينات.

\end{RLtext}

\noindent\rule[2pt]{\textwidth}{0.5pt}

\begin{RLtext}

{\textbf{الكلمات المفتاحية:}}

\LR{STM32Cube}، قياس الأداء، البصمة الذاكرية، متحكم دقيق، \LR{FreeRTOS}، \LR{USB}، \LR{GD32}، \LR{Texas Instruments}، مقارنة بين المصنّعين.
\\

\end{RLtext}

\noindent\rule[2pt]{\textwidth}{0.5pt}